%! Polynomial Functions and Complex Zeros — Unit 2, Lesson 2.4 — Homework (Answer Key)
\documentclass[10pt]{article}
\usepackage{precalculus-article}
\usepackage{precalculus-key}   % pulls in -boxes; do NOT also load -boxes
\newcommand{\TallMath}[1]{$\displaystyle #1\rule[-1.4em]{0pt}{3.2em}$}

\begin{document}

\pageheader{Unit 2, Lesson 2.4}{Homework --- Answer Key}
\namedateperiod

\vspace{0.1in}

\begin{enumerate}[leftmargin=1.5em, itemsep=12pt]

\item Write a polynomial with real coefficients of \emph{minimum degree} that has zeros
      at $x = 3$ (multiplicity 2) and $x = 1 - 2i$. Leave your answer in factored form.
      State the degree.

      \vspace{0.06in}
      $p(x) = $ \ans{$(x-3)^2\bigl(x-(1-2i)\bigr)\bigl(x-(1+2i)\bigr)$}
      \hfill Degree: \ans{4}

      \begin{teachernote}
      Expanded: $(x-3)^2\bigl((x-1)^2+4\bigr)=(x-3)^2(x^2-2x+5)$. Accept either form.
      \end{teachernote}

\item A degree-5 polynomial with real coefficients has zeros at $x = -2$,
      $x = 0$ (multiplicity 2), and $x = 1 + i$.
      List \emph{all five} zeros.

      \ansline{$x=-2$ (mult.\ 1), $x=0$ (mult.\ 2), $x=1+i$, $x=1-i$. Total: $1+2+1+1=5$. \checkmark}

\item Let $p(x) = x^4 - 6x^2 + 9$.
      \begin{enumerate}[leftmargin=1.5em, itemsep=6pt]
        \item Factor $p(x)$ completely.
              \ansline{$p(x)=(x^2-3)^2=(x-\sqrt{3})^2(x+\sqrt{3})^2$}
        \item List all zeros and their multiplicities.
              \ansline{$x=\sqrt{3}$ (multiplicity 2) and $x=-\sqrt{3}$ (multiplicity 2).}
        \item Crosses or tangent at each zero?
              \ansline{Tangent (touches, bounces) at both zeros — both have even multiplicity (2).}
      \end{enumerate}

\item Is $g(x) = 2x^3 - 5x$ even, odd, or neither?

      \vspace{0.06in}
      $g(-x) = $ \ans{$2(-x)^3 - 5(-x)$}

      \vspace{0.08in}
      Simplified: \ans{$-2x^3 + 5x = -(2x^3 - 5x) = -g(x)$}

      \vspace{0.08in}
      $g$ is (circle): \textbf{Even} \quad {\color{keyred}\textbf{Odd}} \quad \textbf{Neither}

\item \textbf{AP-Style Multiple Choice.}
      $p(x) = (x+3)^2(x-1)(x-5)^3$.

      \begin{enumerate}[leftmargin=1.5em, itemsep=3pt, label=(\Alph*)]
        \item {\color{keyred}\textbf{(A) $x=-3$ (tangent), $x=1$ (crosses), $x=5$ (crosses)
              $\leftarrow$ correct}} \\
              Multiplicity 2 at $x=-3$ (even $\Rightarrow$ tangent); mult.\ 1 at $x=1$
              (odd $\Rightarrow$ crosses); mult.\ 3 at $x=5$ (odd $\Rightarrow$ crosses).
        \item (B) Incorrect — reverses behavior at $x=-3$ and $x=1$.
        \item (C) Incorrect — mult.\ 3 is odd, so $x=5$ crosses, not tangent.
        \item (D) Incorrect — wrong zeros (sign errors).
      \end{enumerate}

\end{enumerate}

\vspace{0.1in}

\begin{extensionbox}
A polynomial $p$ of unknown degree has output values recorded at equally-spaced inputs.
The third differences of those outputs are all constant, each equal to $6$.
\begin{enumerate}[leftmargin=1.5em, itemsep=5pt]
  \item What is the degree of $p$?
        \ansline{Degree 3. The $n$th differences of a degree-$n$ polynomial are constant (and nonzero). The 3rd differences are constant, so the degree is 3.}
  \item Write a general form of a polynomial of that degree:
        \ans{$p(x)=ax^3+bx^2+cx+d$} (with $a\neq 0$)
\end{enumerate}
\end{extensionbox}

\vspace{0.1in}

\begin{reflectionbox}
\textbf{Preview:} In Lesson 1.6 we will study rational functions. As you finish this
homework, think about: if a polynomial can be written as a ratio $\dfrac{p(x)}{q(x)}$,
what happens when $q(x) = 0$? We'll explore vertical asymptotes and holes next class.
\end{reflectionbox}

\end{document}
